\documentclass[10pt,a4paper]{article}
\usepackage{amsmath,amssymb,amsthm,graphicx}
\usepackage{xfrac}
\newenvironment{enumeratei}
{\bgroup\def\theenumi{\roman{enumi}}\def\theenumii{\arabic{enumii}}\begin{enumerate}}
{\end{enumerate}\egroup}

\hyphenation{addi-tion-al belong belongs denote exchan-ge exchan-ged exchan-ges}

\newcommand{\psfrac}[2]{\sfrac{(#1)}{#2}}

\newcommand\mto{\mathchoice{\longmapsto}{\mapsto}{\mapsto}{\mapsto}}
\let\oldbigsqcup\bigsqcup
\renewcommand{\bigsqcup}{\mathop{\textstyle\oldbigsqcup}\displaylimits}

\let\epsilon\varepsilon

\DeclareMathOperator{\Ess}{Ess}
\DeclareMathOperator{\Leb}{Leb}

\usepackage{tikz}
\usetikzlibrary{plotmarks}

\theoremstyle{plain}
\newtheorem{theorem}{Theorem}[section]
\newtheorem{lemma}[theorem]{Lemma}
\newtheorem{proposition}[theorem]{Proposition}
\newtheorem{corollary}[theorem]{Corollary}
\newtheorem{claim}{Claim}

\theoremstyle{definition}
\newtheorem{remark}[theorem]{Remark}
\newtheorem{example}[theorem]{Example}

\newcommand{\A}{\mathcal{A}}
\newcommand{\B}{\mathcal{B}}
\newcommand{\I}{\mathcal{I}}

\newcommand{\Z}{\mathbb{Z}}
\newcommand{\N}{\mathbb{N}}
\newcommand{\R}{\mathbb{R}}
\title{A centre whose orbit meets the discontinuities lies in a non-trivial connection}
\author{Przemyslaw Berk, Frank Trujillo}
\date{doi:10.5802/jep.302; CC BY4.0}
\begin{document}
\maketitle
\noindent\textit{Editorial note.} The counted statement is the author's labelled lemma environment at source0.tex lines 652-670; its proof environment opens at line 672, the line after the blank separator, so the pairing is by adjacency and the retained proof range 672-682 is the whole environment.\par\medskip
\noindent\textit{Editorial note.} Four numbered environments are retained as context because the statement or the proof refers to them: eq:sym\_identity\_endpoints\_2 (source lines 480-482), eq:sym\_identity\_endpoints (source lines 475-478), eq:inverse\_iterates (source lines 533-537) and eq:delta\_1/2 (source lines 539-544).\par\medskip
\noindent\textit{Editorial note.} This source is new to the material-authoring work and its preamble is built from the author's own macro and theorem block (source lines 8-45) rather than from a prepared adapter.\par\medskip
\noindent\textit{Editorial note.} This source has no label-number dump and the retained counter values are left empty, so the printed numbering of the card is the wrapper's default; the identity of the counted item is its author label.\par\medskip
\noindent\textit{Editorial scope.} One original result of the pinned publisher article, counted once: a centre whose orbit meets the discontinuities lies in a non-trivial connection (author label lem:symmetric\_connection).\par\medskip
\noindent\textit{Allowed standard prerequisites (editorial).}\begin{itemize}
\item Symmetry and the conjugacy equation. The reflection of the interval conjugates the transformation to its inverse off the discontinuities.. It is the hypothesis and the source of the endpoint identities.
\item Centre of a continuity interval. c\_beta is the middle point of the continuity interval indexed by beta, and c\_\{1/2\} is the middle of the whole interval.. It is the point whose orbit is followed.
\item Connection. A connection is a finite orbit fragment joining two discontinuities.. It is the structure that the conclusion places the centre in.
\end{itemize}\medskip
\begin{equation}\label{eq:sym_identity_endpoints_2}
T^{-1}\circ \I_I(\partial I_{\alpha})=\partial I_{\hat\alpha},\ \text{where}\ \pi_0(\hat\alpha)=\pi_0(\alpha)-1,
\end{equation}

\begin{equation}
\label{eq:sym_identity_endpoints}
\I_I\circ T(\partial I_{\alpha})=\partial I_{\hat\alpha},\ \text{where}\ \pi_0(\hat\alpha)=\pi_0(\alpha)+1,
\end{equation}

\begin{equation}
\label{eq:inverse_iterates}
T^{-m}(c_\alpha) =
{T^{-1} \circ \I_I} \big( T^{m - \delta_{1/2}(\alpha)}(c_\alpha) \big),
\end{equation}

\begin{equation}
\label{eq:delta_1/2}
\delta_{1/2}(\beta) = \left\{\begin{array}{ll} 1 & \text{ if } \quad \beta = \sfrac{1}{2}, \\
0 & \text{ otherwise}.
\end{array} \right.
\end{equation}

\begin{lemma}
\label{lem:symmetric_connection}
Let $T: I \to I$ be a symmetric IET and $\beta \in \A \cup \{\sfrac{1}{2}\}$. Assume there exists $m \in \Z$ and $\alpha \in \A$ such that $T^m(c_\beta) = \partial I_\alpha$ and $T^k(c_\beta) \notin \{\partial I_\gamma\}_{\gamma \in \A}$, for any $-|m| < k < |m|$.

Then {if $m\ge 0$}, we~have $\alpha \neq \pi_0^{-1}(1)$ and\vspace*{-3pt}
\[
T^{-m - \delta_{1/2}(\beta)}(c_\beta) = \partial I_{\hat \alpha},
\]
where $\delta_{1/2}$ is given by \eqref{eq:delta_1/2} and $\pi_0( \hat \alpha) = \pi_0(\alpha) - 1$.

{If on the other hand $m<0$, then $\alpha\neq \pi_1^{-1}(1)$ and\vspace*{-3pt}
\[
T^{-m - \delta_{1/2}(\beta)}(c_\beta) = \partial I_{\hat \alpha},
\]
with $\pi_0( \hat \alpha) = \pi_0(\alpha) + 1$.
}

In particular, $c_\beta$ lies inside a non-trivial connection.
\end{lemma}

\begin{proof}

If $m = 0$, then necessarily $\beta = \sfrac{1}{2}$ and we have $c_{1/2} = \partial I_\alpha$ with $\pi_1(\alpha) \neq 1$. By \eqref{eq:sym_identity_endpoints_2}, $T^{-1}(c_{1/2}) = \partial I_{\hat \alpha}$, where $\pi_0( \hat \alpha) = \pi_0(\alpha) - 1$.

Without loss of generality, let us assume $m < 0$, the case $m > 0$ being analogous. By \eqref{eq:inverse_iterates},
\begin{equation*}
\label{eq:inverse_iterates_negative}
T^m(c_{\beta})=T^{-1}\circ \I_I(T^{-m-\delta_{1/2}(\beta)})\iff T^{-m- \delta_{1/2}(\beta) }(c_\beta) = \I_I\circ T \big( T^{m}(c_\beta)\big).
\end{equation*}
Notice that since $T^{m + 1}(c_\beta) \notin \{\partial I_\gamma\}_{\gamma \in \A}$ then $T^{m}(c_\beta) = \partial I_\alpha \neq \partial I_{\pi_0^{-1}(1)}$. Hence, by~\eqref{eq:sym_identity_endpoints}, the equation above implies $T^{-m - \delta_{1/2}(\beta)}(c_\beta) = \partial I_{\hat \alpha}$, where \hbox{$\pi_0( \hat \alpha) = \pi_0(\alpha) + 1$}.
\end{proof}

\end{document}
